\subsubsection{Impedir acceso a cuenta en estado pendiente}

\textbf{Descripción:} Si el usuario aún no fue aprobado, el login devuelve mensaje de bloqueo y no genera token.

\vspace{0.5em}
\noindent\textbf{PreCondiciones:}
\begin{itemize}[nosep]
    \item Usuario creado con estado pendiente (rol null/activo false) o bandera interna.
\end{itemize}

\vspace{0.5em}
\noindent\textbf{PostCondiciones:}
\begin{itemize}[nosep]
    \item Sin sesión creada; se informa motivo al cliente.
\end{itemize}

\vspace{0.5em}
\noindent\textbf{Actores:} Solicitante (usuario recién registrado).

\vspace{0.5em}
\noindent\textbf{Flujo de eventos principal:}
\begin{enumerate}[nosep]
    \item Usuario intenta login.
    \item Backend consulta estado y detecta ``PENDIENTE APROBACIÓN''.
    \item Lanza excepción identificada (``USUARIO PENDIENTE APROBACIÓN'').
    \item Controlador traduce a mensaje legible.
    \item Usuario permanece sin acceso hasta aprobación.
\end{enumerate}

\vspace{0.5em}
\noindent\textbf{Flujos alternativos:}

\vspace{0.3em}
\noindent\textit{A1. Usuario ya aprobado}
\begin{itemize}[nosep]
    \item Flujo normal de login (con o sin 2FA).
\end{itemize}

\vspace{0.3em}
\noindent\textit{A2. Credenciales inválidas}
\begin{itemize}[nosep]
    \item Error de autenticación estándar.
\end{itemize}

\vspace{0.5em}
\noindent\textbf{Requerimientos especiales:}
\begin{itemize}[nosep]
    \item Mensaje específico en AuthController.
    \item SeguridadService lanza excepción para estado pendiente.
\end{itemize}
